\chapter{一致性模型与连贯性协议的定义与验证}

\begin{epigraph}
\textbf{核心命题}　如何确保一个一致性模型或连贯性协议是"对"的？本章讨论操作化定义与公理化定义两类形式化方式、行为探索（litmus 测试）、以及形式化方法（模型检测、定理证明）与硬件测试两类验证手段。
\end{epigraph}



\section{定义}

一致性模型与连贯性协议都需要严格的\textbf{定义}——否则"对错"无从谈起。有两类风格化定义。

\subsection{操作化定义}

\textbf{操作化定义}（operational definition）用一个\textbf{抽象机}（abstract machine）描述合法行为：

\begin{itemize}
\item SC 的"全局交换开关 / 任意时刻仅一核访存"；
\item TSO 的"每核 FIFO 写缓冲 + 一把全局锁"；
\item ARMv8 的 Flowing/operational 模型。
\end{itemize}


优点：直观、易于人工理解；缺点：难以穷举所有行为，难以直接用于自动化验证。

Owens/Sarkar/Sewell 的 x86-TSO（TPHOLs 2009）给出了操作化模型的经典示例。

\subsection{公理化定义}

\textbf{公理化定义}（axiomatic definition）用一组\textbf{公理}刻画合法执行：

\begin{itemize}
\item \textbf{happens-before}（happens-before 关系）；
\item \textbf{preserved program order}（PPO，被保持的程序序）；
\item \textbf{coherence}（同地址的操作顺序）；
\item \textbf{acyclicity}（关系图无环）。
\end{itemize}


优点：适合形式化验证（可直接输入模型检测器/定理证明器）；缺点：抽象，人工理解门槛高。

代表工具链：\textbf{cat} 语言（DSL，描述公理）+ \textbf{herd} 模拟器（在 cat 模型上模拟 litmus 测试）。Alglave 等在《Herding Cats》（\emph{ACM TOPLAS} 2014，被引超 536）中提出这一统一框架，对 x86、ARM、POWER 给出统一的 cat 模型，并建立 litmus 标准命名（mp/sb/lb/wrc/rwc 等）。

\section{探索内存一致性模型的行为}

\subsection{基准测试}

\textbf{Litmus 测试}（litmus test）是探索内存模型行为的标准工具——小段并发代码 + 预期结果，用于区分内存模型。经典 litmus：

\begin{itemize}
\item \textbf{Store Buffering（SB）}：两线程各写对方将读的位置后再读对方位置；\texttt{r1=0, r2=0} 的结果\textbf{在 SC 下被禁止，在 TSO 下被允许}（源于每核 store buffer）。这是区分 TSO 与 SC 的经典测试。
\item \textbf{Message Passing（MP）}：在 SC 与 TSO 下都被禁止，但在 ARM/POWER 无 fence 的松弛模型下被允许——区分 TSO 与松弛模型的经典测试。
\end{itemize}


\subsection{探索}

\textbf{herdtools7}（Alglave \& Maranget，Inria）是探索内存模型行为的标准工具链：

\begin{itemize}
\item \textbf{diy7}：自动生成 litmus 测试；
\item \textbf{herd7}：在 .cat 模型上模拟结果；
\item \textbf{litmus7}：在真实硬件上随机化运行数百万次，做 hardware probe。
\end{itemize}


近期 GPU 专用工具：Dartagnan（Ponce de Leon 等，形式化建模）与 GPUHarbor（Sorensen \& Levine，基于 WebGPU 的浏览器化 GPU 内存测试平台）。

\section{验证实现}

定义清楚后，需要\textbf{验证实现}符合定义。两类手段。

\subsection{形式化方法}

\textbf{形式化方法}（formal methods）用数学手段证明实现的正确性：

\begin{itemize}
\item \textbf{模型检测}（model checking）：显式状态枚举（Murphi，David Dill, Stanford）或符号模型检测（SMV/BDD）。广泛用于 cache coherence 协议验证（CMU 课程实验即用 Murphi）。FLASH 协议的参数化验证（CHARME 2001）是经典案例。
\item \textbf{定理证明}（theorem proving）：Coq、Isabelle/HOL。x86-TSO 即在 HOL 中机械化（Owens/Sarkar/Sewell）。
\item \textbf{TLA+}：用于协议规约与验证，工业界（如 Intel、Amazon）广泛使用。
\end{itemize}


\subsection{测试}

\textbf{测试}（testing）在真实硬件上捕捉违规行为：

\begin{itemize}
\item litmus7 的随机化运行（数百万次迭代）；
\item 硬件探针。
\end{itemize}


测试不能证明正确性，但能发现 bug——尤其对于商用 ISA 文档不全或实现偏离规范的情况。

\section{历史与进一步阅读资料}

\subsection{Alpha 21164 的教训}

\textbf{DEC Alpha 21164} 因不保证数据依赖序（data-dependency ordering）而著名：一个读可能在挂起的连贯性失效/写提交完成前返回\textbf{本地缓存的陈旧值}，效果上"读到未提交/陈旧数据"。这破坏了 double-checked locking 等惰性初始化模式，迫使 Linux 内核历史上在 Alpha 上使用显式屏障（\texttt{smp\_read\_barrier\_depends()}）。这是弱内存模型踩坑的教科书案例。

\subsection{x86 "processor ordering" 的澄清}

Intel 在 2007–2008 年的澄清（Intel SDM rev.27 / AMD 文档 3.17，2008 年 11 月）把早期模糊的"processor ordering"收敛为 TSO 形式化。白皮书《Intel 64 Architecture Memory Ordering White Paper》（文档号 318147，2007）是这一澄清的关键文件。Sarkar/Sewell 等的 Cambridge 团队推动了这一形式化。

\subsection{ARMv8 multicopy-atomic 修订}

ARMv8 于 2017 年修订为 multicopy-atomic 模型（旧模型允许"写先对部分线程可见再对所有线程可见"，但因未被产品实现利用而去除）。Pulte 等在 POPL 2018《Simplifying ARM Concurrency》（被引超 240）给出 axiomatic 与 operational 两套等价模型。

\section{小结}

一致性模型与连贯性协议的定义有两类风格：操作化（抽象机，直观）与公理化（happens-before + acyclicity，适合验证）。Litmus 测试（SB 区分 SC/TSO，MP 区分 TSO/松弛）是探索模型行为的标准工具。herdtools7（diy7/herd7/litmus7）是 litmus 工具链。验证手段分形式化方法（Murphi 模型检测、Coq/Isabelle 定理证明、TLA+）与硬件测试。历史教训：Alpha 21164 的数据依赖漏洞、x86 processor ordering 的 TSO 澄清（2007–2008）、ARMv8 的 multicopy-atomic 修订（2017）。

\section*{参考文献}
\begin{itemize}
\item Owens S., Sarkar S., Sewell P. A better x86 memory model: x86-TSO. TPHOLs 2009. https://www.cl.cam.ac.uk/\textasciitilde{}pes20/weakmemory/x86tso-paper.tphols.pdf
\item Alglave J., Maranget L., Tautschnig M. Herding cats. \emph{ACM TOPLAS}, 2014. http://www0.cs.ucl.ac.uk/staff/j.alglave/papers/toplas14.pdf
\item herdtools7. https://github.com/herd/herdtools7
\item Pulte C., Flur S., Deacon W., et al. Simplifying ARM concurrency: Multicopy-atomic axiomatic and operational models for ARMv8. POPL 2018. https://doi.org/10.1145/3158107
\item Intel. \emph{Intel 64 Architecture Memory Ordering White Paper}. 318147, 2007. https://www.cs.cmu.edu/\textasciitilde{}410-f10/doc/Intel\_Reordering\_318147.pdf
\item RISC-V. \emph{RVWMO Explanatory Material}. https://docs.riscv.org/reference/isa/v20260120/unpriv/mm-eplan.html
\item Russ Cox. \emph{Hardware memory models}. https://research.swtch.com/hwmm
\item Bornholt J. \emph{Memory models tutorial}. https://jamesbornholt.com/blog/memory-models/
\item Lustig D., et al. Automated synthesis of comprehensive memory model litmus test suites. ASPLOS 2017. https://research.nvidia.com/sites/default/files/pubs/2017-04\_Automated-Synthesis-of/ASPLOS\_2017\_Memory\_Model\_Verification.pdf
\item CMU. \emph{Murphi coherence lab}. https://users.ece.cmu.edu/\textasciitilde{}bgold/teaching/coherence.html
\item Flash parameterized verification. CHARME 2001. http://mcmil.net/pubs/CHARME01.pdf
\item UMD. \emph{Alpha reordering}. https://www.cs.umd.edu/\textasciitilde{}pugh/java/memoryModel/AlphaReordering.html
\item McKenney P. \emph{Memory ordering in modern microprocessors}. http://www.rdrop.com/users/paulmck/scalability/paper/ordering.2007.09.19a.pdf
\item Cambridge weak memory project. https://www.cl.cam.ac.uk/\textasciitilde{}pes20/weakmemory/
\item SIGARCH blog. \emph{GPU memory consistency testing}. https://www.sigarch.org/gpu-memory-consistency-specifications-testing-and-opportunities-for-performance-toolting/
\item Sorin D. J., Hill M. D., Wood D. A. \emph{A Primer on Memory Consistency and Cache Coherence} (2nd ed.). 2020, Chapter 11. https://pages.cs.wisc.edu/\textasciitilde{}markhill/papers/primer2020\_2nd\_edition.pdf
\end{itemize}
